\begin{lemma}
For every positive integer $k$ there is $D_0$ such that the following holds
for all $D\ge D_0$.  Let $\mathcal F\in H(k,k,D)$, and let
$f\in E(\mathcal F)$ maximize $b_{L(\mathcal H),kD}(e)$ over all
$\mathcal H\in H(k,k,D)$ and all $e\in E(\mathcal H)$.  Then every vertex of
$f$ has degree $D$, and $f$ intersects exactly $k(D-1)$ other edges of
$\mathcal F$.
\end{lemma}
\begin{proof}
This is the full $k$-uniform, $k$-simple form of the saturation step used as
Lemma~\ref{lem: structure1} in the paper.  We spell out the perturbation
argument because it is later used through this node, not as an uncited paper
fact.
Take $D_0\ge 2$.  Then every $D\ge D_0$ is large enough for the fresh
vertices introduced below, whose degrees are at most two, to respect the
maximum-degree bound.

By \noderef{HypergraphClass}, the assumptions $\mathcal F\in H(k,k,D)$ mean
that $\mathcal F$ is $k$-uniform, $k$-simple, and has maximum degree at most
$D$.  The first condition is \noderef{UniformHypergraph}; the second is
\noderef{TSimpleHypergraph}; and the degree condition is
\noderef{MaxDegreeAtMost}, with vertex degrees interpreted by
\noderef{HypergraphDegree}.  The line graph is \noderef{LineGraphOfHypergraph},
and the local parameter is
\[
b_{L(\mathcal F),kD}(f)
=\frac{\deg_L(f)}{kD}
-\frac{P_L(f)}{(kD)^2}
+\frac{T_L(f)}{(kD)^3}
\]
by \noderef{LocalBParameter}, where $P_L$ and $T_L$ are the independent-pair
and independent-triple counts of \noderef{IndependentPairCount} and
\noderef{IndependentTripleCount}.

First suppose that a vertex $v\in f$ has degree less than $D$.  Add one new
edge containing $v$ and $k-1$ fresh vertices, and call the enlarged
hypergraph $\mathcal F_1$.  It is still $k$-uniform because the new edge has
exactly $k$ vertices and all old edges are unchanged.  It is still
$k$-simple because the new edge meets every old edge only possibly at $v$,
so no intersection can exceed $k$.  The maximum degree remains at most $D$:
only the degree of $v$ increases by one, and it was below $D$, while the
fresh vertices have degree one.  Thus
\noderef{HypergraphClass} still places $\mathcal F_1$ in $H(k,k,D)$.

Let $f_1$ be the old edge $f$ inside $\mathcal F_1$.  The new edge is a new
neighbor of $f_1$ in the line graph, so \noderef{LineGraphOfHypergraph}
increases $\deg_L(f)$ by one.  All adjacencies among old neighbors of $f$
are unchanged.  The new independent pairs are exactly pairs consisting of
the new edge and an old neighbor of $f$ that does not contain $v$; the new
independent triples are exactly the new edge together with an old independent
pair whose two members do not contain $v$.  In particular the triangle term
does not decrease, while the possible increase in the independent-pair term
is at most the old line-graph degree of $f$.  Since
\noderef{MaxDegreeAtMost} gives at most $D-1$ other edges through each of
the $k$ vertices of $f$, the old degree is at most $k(D-1)<kD$.  Therefore
the change in \noderef{LocalBParameter} is at least
\[
\frac1{kD}-\frac{kD-1}{(kD)^2}>0.
\]
This contradicts the assumed maximality of $(\mathcal F,f)$.  Hence every
vertex of $f$ has degree exactly $D$.

Now count incidences between vertices of $f$ and edges other than $f$.  Since
every vertex of $f$ has degree $D$, there are exactly $k(D-1)$ such
incidences.  Each edge meeting $f$ contributes at least one of these
incidences, so $f$ intersects at most $k(D-1)$ other edges.  If it
intersected fewer than $k(D-1)$ distinct other edges, then some neighboring
edge $g$ would contain at least two vertices of $f$.

In that case perform the following splitting operation, keeping the local
counts explicit.  Choose distinct vertices $v,u\in g\cap f$.  Introduce
fresh vertices
$x_1,\ldots,x_{k-1}$, replace the edge index $g$ by a modified edge
\[
  g'=(g\setminus\{v\})\cup\{x_1\},
\]
and add a new edge
\[
  a=\{v,x_1,\ldots,x_{k-1}\}.
\]
The edge $g'$ still has $k$ vertices, because one old vertex was removed and
one fresh vertex was inserted, and $a$ has $k$ vertices; hence
\noderef{UniformHypergraph} is preserved.  The new edge $a$ meets old edges
only possibly at $v$, and it meets $g'$ only at $x_1$, so its intersection
with every other edge has size at most one.  For an old edge $h$, the
intersection $g'\cap h$ is obtained from $g\cap h$ by possibly deleting
$v$ and by adding no old vertex, since $x_1$ is fresh; therefore
$|g'\cap h|\le |g\cap h|\le k$.  Thus \noderef{TSimpleHypergraph} is
preserved.  The maximum degree bound is also preserved: the degree of $v$
loses the old incidence with $g$ and gains the incidence with $a$, every
other old vertex has the same or smaller degree, and each fresh vertex has
degree at most two.  Hence the new hypergraph $\mathcal F_2$ still lies in
$H(k,k,D)$ by \noderef{HypergraphClass}.

Let $f_2$ be the copy of $f$ in $\mathcal F_2$, and let
$d=\deg_{L(\mathcal F)}(f)$.  Since $g'$ still contains the vertex
$u\in f$, it remains adjacent to $f_2$ in the line graph
\noderef{LineGraphOfHypergraph}; the new edge $a$ is also adjacent to
$f_2$ through $v$.  No old neighbor other than $g$ is removed.  Thus
\[
  \deg_{L(\mathcal F_2)}(f_2)=d+1 .
\]

Now compare independent triples in the two line-graph neighborhoods, using
\noderef{IndependentTripleCount}.  An old independent triple not containing
$g$ is unchanged.  An old independent triple containing $g$ has the form
$\{g,h_1,h_2\}$ where $h_1$ and $h_2$ are disjoint from $g$ and from each
other.  Because $x_1$ is fresh and $g'$ is obtained from $g$ by deleting
$v$ and adding $x_1$, both $h_1$ and $h_2$ are disjoint from $g'$ as well.
Therefore $\{g',h_1,h_2\}$ is an independent triple in the new
neighborhood.  This gives an injection from old independent triples to new
independent triples, so
\[
  T_{L(\mathcal F_2)}(f_2)\ge T_{L(\mathcal F)}(f).
\]

It remains to bound the increase in independent pairs, as defined in
\noderef{IndependentPairCount}.  First, every old independent pair injects
into a new independent pair: pairs not using $g$ are unchanged, and a pair
$\{g,h\}$ with $g\cap h=\emptyset$ maps to $\{g',h\}$, which is still
independent because $h$ is old and avoids both $g$ and the fresh vertex
$x_1$.  Hence any increase in $P$ consists only of new independent pairs
not in this injection.  There are two possible sources.

The first source is a pair involving the added edge $a$.  Such a pair can
only use an old neighbor $h$ of $f$ that does not contain $v$, because
$a$ contains $v$ and is adjacent to every edge containing $v$; also $a$ is
adjacent to $g'$ through $x_1$.  Since the first part of the proof already
showed $\deg_{\mathcal F}(v)=D$, exactly $D-1$ old neighbors of $f$ contain
$v$ (all edges through $v$ except $f$ itself), one of which is $g$.  Thus
there are $d-(D-1)$ possible independent pairs involving $a$.

The second source is a pair $\{g',h\}$ which was not the image of an old
independent pair.  This means $g$ and $h$ were adjacent before the split but
$g'$ and $h$ are disjoint after it.  Since $g'$ differs from $g$ only by
removing $v$ and adding the fresh vertex $x_1$, this can happen only when
$g\cap h=\{v\}$.  Such an $h$ is one of the old neighbors through $v$ other
than $g$, so there are at most $D-2$ of them.  Consequently
\[
  P_{L(\mathcal F_2)}(f_2)-P_{L(\mathcal F)}(f)
  \le (d-(D-1))+(D-2)=d-1.
\]
Our contradiction assumption says $d<k(D-1)$, so $d-1<kD$.  Therefore, with
$M=kD$, \noderef{LocalBParameter} gives
\[
\begin{aligned}
b_{L(\mathcal F_2),M}(f_2)-b_{L(\mathcal F),M}(f)
&=\frac{1}{M}
  -\frac{P_{L(\mathcal F_2)}(f_2)-P_{L(\mathcal F)}(f)}{M^2}
  +\frac{T_{L(\mathcal F_2)}(f_2)-T_{L(\mathcal F)}(f)}{M^3}  \\
&\ge \frac{1}{M}-\frac{d-1}{M^2}
>0 .
\end{aligned}
\]
This produces another member of $H(k,k,D)$ with strictly larger local
$b$-parameter than $(\mathcal F,f)$, contradicting the assumed maximality.
Therefore no neighboring edge contains two vertices of $f$, and the number
of distinct neighboring edges equals the incidence count $k(D-1)$.
\end{proof}
